\chapter{Acknowledgements}
\thispagestyle{empty}
Questa tesi non esisterebbe se non avessi ricevuto il supporto di tantissime persone. Supporto tecnico da parte di Professori, Dottorandi e Compagni di corso che mi hanno aiutato ad appassionarmi alle materie studiate, a capirle meglio, a ripassarle, e a superare gli esami. Supporto morale da parte della mia Famiglia e degli Amici che mi hanno sostenuto, accompagnato, e fatto crescere.

Desidero, in questa pagina, ringraziare il professor Roversi e la professoressa Bono che sono stati, innanzitutto, degli ottimi insegnati, mi hanno fatto amare le loro materie e i loro corsi. Con la professoressa Bono ho scritto anche la tesi triennale: le sono davvero grato per avermi contattato lei per prima e aver sciolto le mie incertezze. Con entrambi ho invece scritto questa tesi, ho potuto farlo con serenità e con i miei tempi, sentendomi sempre seguito, e sapendo che il mio lavoro era controllato e corretto in tempi brevissimi. Mi sento molto fortunato ad aver trovato un gruppo di lavoro del genere. Gruppo che è costituito anche dai dottorandi, in particolare mi preme ringraziare Mario e Matteo che mi hanno dato preziosissimi consigli durante la scrittura.

Un pensiero non può che andare alla mia famiglia, al Papà e alla Mamma che mi vogliono bene e mi supportano a prescindere da tutto, ai Nonni che mi spronano e credono sempre in me e alla Sorellina che, anche se è in giro per il mondo, so che c'è sempre per me.

Un grosso grosso grazie va poi agli amici, dei quali sono veramente tanto orgoglioso, alcuni veramente di vecchissima data che mi porto nel cuore dalle elementari. Io, Eddi, Emi e Luca siamo cresciuti assieme ormai molto più di quanto non siamo cresciuti da soli. Abbiamo giocato, riso, sbagliato e vissuto assieme esperienze che mi hanno fatto diventare chi sono.  Al liceo, ho conosciuto, tra gli altri, Tom e Fra che sono stati miei compagni di classe, Lore, Marco, Alessandro che sono stati compagni di scuola, e Sara che ho conosciuto a Bardonecchia. Con loro, e tanti tanti altri, ho passato cinque anni bellissimi e ho scoperto più in la che le amicizie, quelle vere, non vengono degradate dal tempo e che basta un messaggio per ritrovarsi come se non fosse passato un giorno. Anche l'università è stata occasione per stringere amicizie e ho trovato persone stupende come Amedeo, Matilde, Asia e Margherita a fisica, e Lorenzo, Angelo e Andrea alla triennale di informatica, devo tantissimo del mio percorso a loro. In magistrale ho avuto la fortuna di incontrare Andrea, senza il quale non avrei passato la maggior parte degli esami e con cui ho condiviso molto più che solo lo studio.

Infine un grande grazie e un grosso bacio va ad Alice che mi è stata sempre vicina, mi ha dato supporto nei momenti difficili e ha sopportato le mie crisi e le mie lagne. Mi dà serenità e mi fa sentire amato, e per questo le sono infinitamente grato. Inoltre se la tesi è vagamente piacevole da leggere tanto del merito è suo.
\begin{flushright}
	Grazie a tutti
\end{flushright} 